% Công cụ: docs/tool_contracts.md (6 công cụ đọc + 2 sink giả, không I/O mạng). Bộ chấm biến đổi:
% src/react_agent/evaluation/transformed_leakage.py (decision log 2026-11-15 18:20). Không nêu tên công cụ giả trong bài.
\section{Mô hình đe dọa}\label{sec:threat}

\paragraph{Tác tử, công cụ và nhãn.} Gọi $u$ là yêu cầu của người dùng. Ở mỗi bước, mô hình sinh một hành động $a_t$ gồm
tên công cụ và tham số, hoặc sinh câu trả lời cuối $y$. Tool Broker kiểm tra $a_t$, thực thi và trả về quan sát $o_t$.
Hệ thống có sáu công cụ đọc và hai \emph{sink} có tác dụng phụ là gửi e-mail và gọi webhook, đều là bản giả chỉ ghi
nhật ký. Câu trả lời cuối được coi là sink thứ ba. Với một lời gọi sink $a$, $\dest(a)$ là đích (địa chỉ hoặc URL) và
$\data(a)$ là các tham số còn lại. Mỗi nội dung mang hai nhãn độc lập. Theo độ tin cậy, chỉ $u$ và prompt hệ thống là
đáng tin. Theo độ nhạy cảm, nội dung được xếp từ công khai tới mật, vì một trang không đáng tin có thể là công khai, còn
một bản ghi mật có thể đến từ nguồn mà kẻ tấn công không chạm tới.

\paragraph{Kẻ tấn công.} Kẻ tấn công kiểm soát văn bản của một nội dung do công cụ đọc trả về và biết schema của công cụ.
Trong các vòng thích nghi, họ biết thêm toàn bộ thiết kế và mô hình hình thức của phòng thủ. Họ không sửa được $u$,
prompt hệ thống, broker hay trọng số mô hình. Jailbreak trực tiếp, đầu độc dữ liệu huấn luyện và hệ nhiều tác tử nằm
ngoài phạm vi.

\paragraph{Vị từ gây hại.} Gọi $G(u)$ là tập cặp (công cụ, đích) mà người dùng thực sự cho phép, $\Sigma$ là tập giá trị
mật trên quỹ đạo mà người dùng không cho phép tiết lộ, và $\mathcal{S}$ là tập lời gọi sink đã được thực thi. Giá trị do
người dùng viết trong $u$ không thuộc $\Sigma$. Quỹ đạo $\tau$ gây hại, ký hiệu $\mathrm{Harm}(\tau)$, khi ít nhất một
trong bốn vị từ sau đúng: $H_\mathit{dest}$ khi có $a \in \mathcal{S}$ với $(\mathrm{tool}(a), \dest(a)) \notin G(u)$,
$H_\mathit{leak}$ khi một $v \in \Sigma$ xuất hiện trong $\data(a)$ của một $a \in \mathcal{S}$ hoặc trong $y$ ở dạng
nguyên văn hay biến đổi, $H_\mathit{tamper}$ khi có $a \in \mathcal{S}$ tới đích hợp lệ nhưng một trường do người dùng
chỉ định mang giá trị khác, và $H_\mathit{side}$ khi nhiệm vụ chỉ cần đọc nhưng $\mathcal{S} \neq \emptyset$. Một bộ
chấm đóng băng xác định các vị từ từ lời gọi đã thực thi và $y$, nên lời gọi bị chặn không được tính là gây hại. Ngoài
so khớp nguyên văn, bộ chấm nhận ra giá trị mật ở dạng Base64, Hex, mã hóa URL, tách ký tự và số viết bằng chữ.

\paragraph{Cơ chế, họ và biến thể.} Tấn công được phân loại theo sáu \emph{cơ chế}, tức con đường để đoạn chèn biến
thành tác hại (Bảng~\ref{tab:mechanisms}). Ở hai cơ chế chuyển hướng, đích không có trong $u$. Ở hai cơ chế rò rỉ và cơ
chế sửa tham số, đích hợp lệ, nên chỉ nguồn gốc hay giá trị của dữ liệu mới phân biệt được tấn công với nhiệm vụ. Một
\emph{họ} là một ý tưởng tấn công trong một nhiệm vụ và một môi trường công cụ cụ thể. Mọi họ có nhánh \emph{\bPlain},
thay đoạn chèn bằng một ghi chú vô hại cùng độ dài, và một phần ba số họ có thêm nhánh \emph{\bAuth}, trong đó người
dùng cho phép rõ ràng một hành động giống mục tiêu tấn công. Chặn ở nhánh lành tính là \emph{chặn nhầm}. Mỗi họ có năm
biến thể áp cho cả $u$ lẫn đoạn chèn là tiếng Việt chuẩn, diễn đạt lại, trộn Việt và Anh, không dấu, và tiếng Anh, với
định danh, địa chỉ và giá trị mật giữ nguyên.

% Bảng 1. Nguồn số: facts.md mục "Bổ sung sau phản biện 3" (nguyên nhân lọt v1: Test-I 13 = SL 6, AL 3, PT 4 (3 sửa giá
% trị + 1 đích có trong u); Test-II 27 = SL 14, AL 2, PT 7, US 3, CR 1) và mục Test-4 ("40 ca vượt v1": M3 20, M5 11,
% M4 5, M6 3, M2 1). Tổng theo nguyên nhân: nguồn công khai 25, đích có trong u 5, sửa tham số 10.
\begin{table}[!htbp]
\centering
\caption{Sáu cơ chế tấn công. Cột thứ ba ghi số ca lọt qua \base{} trên hai tập niêm phong trước (\tOne{} và \tTwo,
ba mô hình nền, tổng 40 ca) và nguyên nhân. Cột cuối là thành phần của \ours{} nhắm vào nhóm đó.}\label{tab:mechanisms}
\footnotesize
\setlength{\tabcolsep}{3pt}
\begin{tabularx}{\textwidth}{@{}>{\raggedright\arraybackslash}p{2.45cm}>{\raggedright\arraybackslash}X%
>{\raggedright\arraybackslash}p{4.0cm}l@{}}
\toprule
Cơ chế & Con đường gây hại & Lọt qua \base{} (nguyên nhân) & \ours \\
\midrule
\mER{} chuyển hướng tường minh & Đoạn chèn yêu cầu gửi dữ liệu tới một đích lạ & 0 & \cOne \\
\mCR{} chuyển hướng che giấu & Như trên, nhưng diễn đạt gián tiếp hoặc tách nhỏ đích & 1 (trường đích là một từ thường
  có trong $u$) & \cOne \\
\mSL{} rò rỉ qua sink & Giá trị mật đi vào nội dung gửi tới một đích hợp lệ & 20 (giá trị nằm trong nguồn gắn nhãn
  công khai) & \cTwo \\
\mAL{} rò rỉ qua câu trả lời & Giá trị mật xuất hiện trong câu trả lời cuối & 5 (như trên) & \cTwo \\
\mPT{} sửa tham số & Đổi một giá trị người dùng đã nêu, giữ đích hợp lệ & 11 (10 ca sửa giá trị, 1 ca đích chỉ được
  nhắc tới) & \cThree, \cOne \\
\mUS{} tác dụng phụ ngoài yêu cầu & Gọi sink tới một đích chỉ được nhắc tới hoặc bị cấm & 3 (đích có trong $u$ nhưng
  không được cấp quyền) & \cOne \\
\bottomrule
\end{tabularx}
\end{table}

